%% One paragraph, self-contained. No cross-references, section numbers,
%% structural labels or citations.
%% Every numerical value is emitted from the hash-bound evidence pass.

Directional disturbance information can change both which aircraft receives a
waypoint and how it flies there. We map and quantify a reproducible computational
operating envelope for this information-to-completion interface in a
multi-quadrotor allocation--control stack. The system is a Hungarian
allocation--control stack in a rigid-body simulator: a wind-agnostic arm uses
Euclidean assignment and the stock position controller, while a wind-aware arm
adds an upwind assignment cost and target-position feedforward. Shared seeds,
paired contrasts and a seven-level force ladder define the primary envelope, and
estimate-noise, temporal, pathway, gain, geometry, swarm-size, capture-radius
and feasibility campaigns extend the same endpoint. The envelope has a sharp
transition at \KneeLevel{}~N (\KneePctHover{}~percent of hover thrust), where
completion rises from \KneeAgnostic{} to \KneeAware{} (paired gain
+\KneeDelta{}, surviving the pre-specified signed-rank family correction). Relative
noise through \BestSigma{} preserves the transition, and causal estimates retain
positive gains under constant and gusting profiles. A complete
assignment-by-feedforward factorial attributes the gain to target feedforward,
while paired gains stay positive from \ScaleMinSize{} through \ScaleMaxSize{}
airframes. Rerunning the ladder across capture radii keeps a transition of
comparable size but moves it from \RadiusTightPeakForce{} to
\RadiusWidePeakForce{}~N, so force describes an operating point rather than an
interface constant. The study therefore contributes a simulator operating
envelope, a transparent computational validation contract, and a parameterized
transfer protocol for hardware and field validation; clearance and collision
quantities are retained as benchmark-level feasibility diagnostics rather than
as deployment claims.
